% Kevin. 150 a 250 palabras EN INGLES + maximo 6 palabras clave.
\begin{abstract}
Real-time noise reduction on embedded hardware forces a choice between filtering
strategies whose relative merit depends on how the interference overlaps the
signal of interest. This work reports a three-node distributed system that makes
that choice explicit and measurable. A first node samples the contaminated signal
at 16~kHz, computes a 1024-point FFT and classifies the interference as tonal,
broadband or band-limited from its spectral flatness and peak-to-median ratio.
The diagnosis routes each block to one of two reduction paths: a cascaded IIR
filter operating in the time domain, or a spectral attenuator that modifies FFT
magnitudes while preserving Hermitian symmetry and reconstructs the signal by
50\,\% overlap-add. The nodes exchange blocks over a framed serial link with
sequence numbering and CRC-16/CCITT error detection, which costs nine bytes per
block and allows lost or corrupted blocks to be counted instead of being silently
averaged into the reported metrics. Acceptance criteria were fixed before
measuring, with per-test thresholds derived from desktop experiments in which the
time-domain filter reached 27.5~dB of SNR improvement when the interference fell
outside its passband but none at all when it did not, while the spectral method
remained uniformly moderate at the cost of 13.3\,\% of the useful energy. Five
test cases, three repetitions each, are evaluated on mean squared error, SNR
improvement, useful-energy conservation, interference attenuation and end-to-end
latency, and each result is reported with its variability across repetitions.
\end{abstract}

\begin{IEEEkeywords}
Adaptive noise reduction, fast Fourier transform, IIR filters, overlap-add,
spectral subtraction, embedded systems.
\end{IEEEkeywords}
